% Draft section for the baseline pipeline comparison.
% Intended for insertion into Final_Report.tex under the evaluation/results section.
% Do NOT wrap in \documentclass / \begin{document}.

\subsection{Baseline Pipeline Comparison}
\label{sec:baseline_pipeline}

We compare the graph-based detector against two unsupervised tabular baselines: One-Class SVM and Isolation Forest. The baseline pipeline is implemented in \texttt{src/baselines} and runs from the same cached graph-pipeline outputs used by the rest of the project. Each baseline invocation writes a run-owned directory under \texttt{baseline\_results/<timestamp>/}, including per-variant JSON files, edge scores, \texttt{summary.json}, \texttt{comparison\_table.md}, and a \texttt{figures/} subdirectory. This mirrors the normal graph pipeline layout, where each \texttt{results/<run>/} directory owns its generated figures.

\subsubsection{Protocol}

All methods use the same LANL-2015 edge feature tables and the same three variants: combined authentication-plus-flow features, authentication-only features, and flow-only features. Variant membership and feature whitelists come from the canonical \texttt{src.variants} module, and binary feature handling comes from \texttt{src.types}. Non-binary columns are converted to percentile ranks; binary columns are left unchanged. Edges marked as self-loops or user metadata edges are excluded from evaluation.

Ground truth comes from LANL red-team source-destination pairs. A graph edge is positive when its \texttt{(src, dst)} pair appears in that set. The one-class baselines train only on normal edges, then evaluate on held-out normal edges plus all positive red-team edges. This is the right contract for unsupervised anomaly detection: learn normal behavior first, then score what looks abnormal.

One-Class SVM uses an RBF kernel with \texttt{nu=0.1}. Features are standardized with \texttt{StandardScaler} fitted on the normal training split, and anomaly scores are obtained by negating \texttt{decision\_function()}. Isolation Forest uses 100 trees, \texttt{contamination="auto"}, and seed 42; anomaly scores are obtained by negating \texttt{score\_samples()}.

Thresholds and pair-level metrics are not reimplemented in the baseline package. Baseline score arrays are adapted into the canonical \texttt{DetectionParams} shape and evaluated through \texttt{src.detection.optimize\_threshold} and \texttt{src.detection.compute\_pair\_metrics}. This keeps threshold semantics, recall, false-positive rate, precision, F1, and AUC aligned with the graph detector.

\subsubsection{Results}

Table~\ref{tab:baseline_combined}, Table~\ref{tab:baseline_auth}, and Table~\ref{tab:baseline_flow} report the current baseline run at \texttt{baseline\_results/20260524\_051746}, evaluated against graph-pipeline run \texttt{20260524\_022620}. Only the graph method and the two implemented one-class baselines are reported.

\begin{table}[h]
\centering
\small
\begin{tabular}{lrrrrrr}
\hline
Method & Recall & FPR & F1 & Precision & AUC & Detected pairs \\
\hline
Graph-based & \textbf{0.9448} & 0.0295 & \textbf{0.0371} & 0.0189 & \textbf{0.9756} & \textbf{291/308} \\
One-Class SVM & 0.0162 & \textbf{0.0010} & 0.0204 & \textbf{0.0273} & 0.7678 & 5/308 \\
Isolation Forest & 0.0292 & 0.0050 & 0.0147 & 0.0098 & 0.9388 & 9/308 \\
\hline
\end{tabular}
\caption{Combined authentication-plus-flow baseline comparison. The graph-based method has the strongest recall and AUC; the one-class baselines keep false positives low but miss most red-team pairs.}
\label{tab:baseline_combined}
\end{table}

\begin{table}[h]
\centering
\small
\begin{tabular}{lrrrrrr}
\hline
Method & Recall & FPR & F1 & Precision & AUC & Detected pairs \\
\hline
Graph-based & \textbf{0.9448} & \textbf{0.0293} & \textbf{0.0463} & \textbf{0.0237} & \textbf{0.9781} & \textbf{291/308} \\
One-Class SVM & 0.0130 & 0.0500 & 0.0012 & 0.0006 & 0.5442 & 4/308 \\
Isolation Forest & 0.2695 & 0.0995 & 0.0124 & 0.0063 & 0.9030 & 83/308 \\
\hline
\end{tabular}
\caption{Authentication-only baseline comparison. Authentication events carry most of the signal for lateral movement in this dataset.}
\label{tab:baseline_auth}
\end{table}

\begin{table}[h]
\centering
\small
\begin{tabular}{lrrrrrr}
\hline
Method & Recall & FPR & F1 & Precision & AUC & Detected pairs \\
\hline
Graph-based & 0.0130 & 0.1000 & 0.0006 & 0.0003 & 0.5154 & 4/308 \\
One-Class SVM & \textbf{0.0227} & \textbf{0.0049} & \textbf{0.0219} & \textbf{0.0212} & 0.7500 & \textbf{7/308} \\
Isolation Forest & 0.0162 & \textbf{0.0049} & 0.0157 & 0.0152 & \textbf{0.8246} & 5/308 \\
\hline
\end{tabular}
\caption{Flow-only baseline comparison. Flow-only data is weak across all methods; it is better treated as supporting context than as the main detection signal.}
\label{tab:baseline_flow}
\end{table}

\subsubsection{Interpretation}

The combined and authentication-only variants show the same pattern: the graph-based detector finds nearly all red-team pairs, while the one-class baselines find only a small fraction at their selected thresholds. Isolation Forest has a respectable AUC on the combined and authentication-only variants, which means it can rank some suspicious edges above normal edges, but its operating point still misses most positives. One-Class SVM is more conservative and sometimes has the lowest false-positive rate, but that comes at the cost of very low recall.

The flow-only result is the exception, and it is useful precisely because it is bad. Without authentication semantics, every method struggles. This supports the broader design choice in the pipeline: authentication events are the primary lateral-movement signal, while flow records are useful only as extra context.

The main conclusion is not that the graph pipeline wins every metric. It does not. The graph pipeline wins the metric that matters most for this dataset: red-team pair coverage under the shared pair-level evaluation. The one-class baselines are useful sanity checks, but they score each edge as a tabular point. They do not model the relational context that makes lateral movement visible.

\subsubsection{Reproduction}

The baseline comparison can be regenerated with:

\begin{verbatim}
uv run baselines.py
uv run baselines.py --run_id 20260524_022620
uv run baselines.py --run_id 20260524_022620 --variant combined
\end{verbatim}

Each command writes to \texttt{baseline\_results/<timestamp>/}. The generated figures are stored in \texttt{baseline\_results/<timestamp>/figures/}; the baseline module does not write report-level figure outputs.
